% id: 66-34
% book: 베이직쎈 확률과 통계 · 04 조건부확률
% section: 개념 19 $\mathrm{P}(B)=\mathrm{P}(A\cap B)+\mathrm{P}(\comp{A}\cap B)$를 이용한 확률
% passage: 다음에 답하시오.
% figure: none
% answer: 
% answer_source: 
% variant_level: 0 (원본 전사)
% 이 파일은 build.mjs 가 items.json 에서 만든다. 고칠 때는 items.json 을 고친다.
\smash{\raisebox{1.5mm}{\dmkichul{베이직쎈 확률과 통계 66쪽 34번}}}
3개의 당첨 제비를 포함한 10개의 제비 중에서 $\mathrm{A}$, $\mathrm{B}$ 두 사람이 이 순서대로 임의로 제비를 한 개씩 뽑을 때, $\mathrm{A}$가 당첨 제비를 뽑는 사건을 $A$, $\mathrm{B}$가 당첨 제비를 뽑는 사건을 $B$라 하자. 다음을 구하시오. \dmcondfill{뽑은 제비는 다시 넣지 않는다.}{}
\dmsubprob{1} $\mathrm{B}$가 당첨 제비를 뽑을 확률\dmhint{$\mathrm{P}(B)=\mathrm{P}(A\cap B)+\mathrm{P}(\blank[4em])$\\ $\hphantom{\mathrm{P}(B)}=\mathrm{P}(A)\mathrm{P}(\blank[3.2em])+\mathrm{P}(\comp{A})\mathrm{P}(B\mid \comp{A})$\\ $\hphantom{\mathrm{P}(B)}=\dfrac{3}{10}\cdot\blank+\blank\cdot\dfrac{3}{9}=\blank$}
\dmsubprob{2} $\mathrm{B}$가 당첨 제비를 뽑았을 때, $\mathrm{A}$가 당첨 제비를 뽑았을 확률\dmhint{$\mathrm{P}(A\mid B)=\dfrac{\mathrm{P}(\blank[4em])}{\mathrm{P}(B)}=\dfrac{\blank}{\blank}=\blank$}
